\section{Related Work}
\label{sec:related}
Research into software energy efficiency has primarily focused on two distinct areas: runtime measurement infrastructures and predictive modeling for specific domains. 

\subsection{Software Energy Measurement}

Accurate energy measurement is the prerequisite for optimization. The most reliable method involves external hardware power meters, but their high cost and non-scalability limit their use to lab environments \cite{Hindle2014GreenMining}. Consequently, the research community has shifted toward software-base power models. Intel's RAPL (Running Average Power Limit) has become the de facto standard for accessing on-chip energy counters \cite{Khan2018RAPLinAction}. However, as noted in recent studies, RAPL's update frequency (approx. 1ms) is often too coarse to capture the energy consumption of short-lived code constructs \cite{10.1145/2425248.2425252}. 
To address this granularity issue, tools like ALEA \cite{7429297} and finer-grained profilers for embedded systems has been proposed. These approaches often use probabilistic sampling or instruction-level characterization to attribute energy to specific code regions. 
% {\color{red}While effective for post-hoc analysis, they remain runtime techniques: they require the code to be fully executable and the hardware to be available, preventing their use during the early design or coding phases.}


\subsection{Static Analysis and Energy Prediction}

To move energy assessment earlier in the lifecycle, researchers have explored static analysis and machine learning. 
In the mobile domain, approaches like \textit{MLEE}\cite{ALVI2021100594} and \textit{EcoAndroid} \cite{Ribeiro2021EcoAndroid}, have successfully demonstrated that structural metrics can predict energy consumption for Android applications. Similarly, \textit{FlipFlop} \cite{Sharma2024FlipFlop} utilizes static analysis to optimize GPU kernel configurations for AI workloads. 

For general-purpose languages like Python, tools such as \textit{GreenPy} \cite{Reya2023GreenPyEA} provide application-level energy insights. However, these existing approaches typically operate at the level of entire applications, methods, or specific API calls. They lack the resolution to estimate the energy cost of fundamental control flow blocks (e.g., individual \texttt{for} loops or \texttt{if} conditions) in general-purpose context.  Furthermore, many existing models can rely on deep learning, which can itself be energy-intensive. By contrast, EnCoDe targets this specific granularity gap using lightweight, interpretable models compliant with Green AI principles. 

\begin{comment}
    
Energy-aware software engineering has received increasing attention in recent years, yet existing approaches remain limited in their ability to provide fine-grained, design-time energy feedback to developers. A key limitation stems from the granularity and timing of current measurement infrastructures.

Traditional energy profiling tools typically operate at coarse granularities, measuring consumption at the process, application, or system level. Early work on energy-aware application design emphasized the importance of energy transparency during development, but restricted measurements to process-level attribution \cite{10.1145/1453175.1453180}. More recent surveys of software energy measurement techniques likewise focus on system-level and application-level estimation, without addressing fine-grained attribution at the level of code constructs \cite{castor2024estimatingenergyfootprintsoftware}. As a result, these approaches provide limited guidance for identifying energy-intensive regions within source code.

Recognizing the limitations of coarse-grained profiling, several researchers have attempted to improve measurement resolution. Techniques based on processor energy counters, such as Intel RAPL, have been proposed to stabilize measurements for short-running code paths, but remain constrained by millisecond-scale sampling granularity \cite{10.1145/2425248.2425252}. Probabilistic approaches such as ALEA further attempt to attribute energy to basic blocks using compiler-level traces and sampling-based estimation, trading deterministic accuracy for scalability \cite{7429297}. More recently, fine-grained measurement has been demonstrated for specific execution domains, such as operator-level energy profiling in deep learning frameworks; however, these approaches remain tightly coupled to particular runtimes and hardware assumptions, limiting their general applicability beyond neural network workloads \cite{10.1145/3680470}.

In parallel, an emerging research direction seeks to estimate energy consumption statically from source code properties, enabling feedback before program execution. MLEE applies machine learning to predict the energy consumption of Java methods in Android applications using traditional code metrics, but operates at method-level granularity and relies on coarse-grained runtime attribution \cite{ALVI2021100594}. Other studies have explored structural relationships between code and energy usage through Abstract Syntax Tree (AST) analysis, showing that syntactic patterns correlate with energy behavior; however, these approaches typically cluster or characterize entire methods rather than providing predictive, construct-level estimators suitable for design-time use \cite{ed0319b79e8749448af332fdb505d8bf}. Rule-based tools such as GreenPy provide energy guidance for Python programs through comparative heuristics, but do not generalize to data-driven predictive modeling across unseen code \cite{Reya2023GreenPyEA}.

In summary, runtime profiling approaches provide empirically grounded measurements but operate at coarse granularity and only after code execution, while static and design-time techniques offer earlier feedback but lack fine-grained empirical validation and structural modeling. \textsc{EnCoDe} bridges this gap through three key contributions: (1) \textbf{PowerLens}, a novel measurement infrastructure that achieves reliable sub-millisecond energy readings for microsecond-scale code blocks by systematically addressing RAPL's temporal limitations through amplification, synchronization, and calibration; (2) \textbf{Feature extraction}, which derives rich AST-based representations capturing complexity, density, diversity, and contextual properties at block granularity; and (3) \textbf{Predictive modeling}, which trains machine learning models on empirically measured block-level energy data to enable accurate design-time estimation for unseen code without execution. By unifying fine-grained measurement with design-time prediction, \textsc{EnCoDe} provides actionable, construct-level energy feedback during software development, addressing limitations present in prior work.
\end{comment}
